\begin{idstatement}{C.3.11}{Proved complete cohomological dichotomy}
\label{st:C.3.11}
Under the hypotheses of C.3.1, a very general smooth fibre
is a two-torus, a primary or secondary Hopf surface of
algebraic dimension zero. The Inoue exclusion is C.3.12.
Both Hopf
entries and the torus entry occur, by C.3.6 and D.2.7.
For the entire smooth locus there are exactly the following
two cohomological alternatives:
\[
\begin{array}{c|c|c}
\text{smooth fibres}&
 (h^0(\OO),h^1(\OO),h^2(\OO))&
 \text{coherent direct images}\\ \hline
\text{two-tori}&(1,2,1)&
 R^2f_*\OO_Y\simeq\det R^1f_*\OO_Y\\
b_1=1,\ b_2=0&(1,1,0)&
 R^1f_*\OO_Y\simeq\OO_{\PP^1},\
 R^2f_*\OO_Y=0 .
\end{array}
\tag{C.PK.9}
\]
The indicated cohomology dimensions hold on every
scheme-theoretic full fibre in the respective branch,
and both direct images satisfy base change there.
In the second branch
$H^q(Y,\OO_Y)=(\C,\C,0,0)$ and
$\Pic^0(Y)=\C\beta$ with the exponential convention
of D.1.9. The first line records a cohomological
alternative rather than a positivity assertion about
the splitting of $R^1f_*\OO_Y$.
\end{idstatement}
\begin{proof}
C.3.7, C.3.9, C.3.10 and C.3.12 remove the bielliptic,
Kodaira and Inoue types from C.3.1; C.3.8 gives the Hopf algebraic
dimension. All smooth fibres are diffeomorphic over the
connected smooth base. In the torus branch they have
$b_1=4$, $b_2=6$, and are K\"ahler by the compact-surface
criterion. Their rational first Chern class and second
Chern number vanish by C.3.1. The surface classification
then makes them two-tori: an effective nonzero canonical
divisor has strictly positive degree against a K\"ahler
form, whereas this rational first Chern class has zero
degree. The minimal K\"ahler surfaces of Kodaira dimension
zero with these Betti numbers are exactly the two-tori.
D.1.12 supplies this branch's full-fibre assertions.

In the Hopf branch every smooth fibre has
$b_1=1$ and $b_2=0$. A nonzero holomorphic two-form
is closed and has strictly positive integral against
its conjugate, so it defines a nonzero de Rham
cohomology class. Thus $b_2=0$ forces $p_g=0$.
Since $\chi(\OO)=0$ by C.3.1, one has $q=1$.
By E4, $\mathscr E=R^1f_*\OO_Y$ is a line bundle,
and its nowhere-zero transgression section $\beta$
from D.1.9 trivializes it. The coherent sheaf
$\mathscr G=R^2f_*\OO_Y$ has rank zero, hence finite
length. Hirzebruch--Riemann--Roch on $Y$ and coherent
Leray give the exact equality
\[
 0=\chi(Y,\OO_Y)
   =\chi(\OO_{\PP^1})-\chi(\mathscr E)+\chi(\mathscr G)
   =1-1+\operatorname{length}(\mathscr G).
\]
Therefore $\mathscr G=0$. Top-degree base change and
the parameter-multiplication exact sequence give
$h^2(\OO_{Y_c})=0$ and $h^1(\OO_{Y_c})=1$ on each
full fibre; E4 gives $h^0(\OO_{Y_c})=1$.
Leray now gives the stated global structure-sheaf
cohomology. Finally $H^1(Y;\Z)=0$ and the exponential
sequence identify $\Pic^0(Y)$ with
$H^1(Y,\OO_Y)=\C\beta$, as asserted.
\end{proof}


% C.3.12 from part_C3_Inoue.tex

